\section{Data and study design}
\label{sec:data}

\subsection{Data, eligibility and selection}
\label{sec:data:sample}

All load series come from the Building Data Genome Project 2 (BDG2) \citep{miller2020bdg2}: hourly electricity readings from university, public and commercial buildings, with site metadata and hourly air temperature. The repository files \citep{bdg2repo} carry a CC BY-SA 4.0 license and the data paper a CC BY 4.0 license, so derived panels keep the share-alike terms. We use the cleaned electricity file for 2016 and 2017. Year 2016 is development data only: thresholds, the daylight-saving mapping and pipeline checks were fixed on it, and no 2016 number is a result. Year 2017 is held out and scored once after the freeze (Section~\ref{sec:data:plan}).

Five eligibility rules use thresholds fixed from 2016. A meter is eligible (E1) with at least 300 valid days in each year (24 readings and a positive sum); this leaves {{sample.n_meters_e1|int}} meters. A day is flat (E2) if four or more consecutive readings between 10:00 and 20:00 are identical. We drop a meter as target and donor (E3) if more than 20 percent of its warm-season weekdays in the scored year are flat, removing {{sample.n_excl_flat_meter|int}} meters. Duplicate meters at a site (correlation of at least 0.999) are dropped once, {{sample.n_dup_dropped|int}} in all. A site enters the sample (E4) if at least 12 meters survive E1 to E3, which excludes Bobcat, Gator and Moose. For each hypothesis a site also needs at least 10 building-events with both compared methods present (E5); exclusions are logged by site and reason. The sample has {{sample.n_buildings|int}} buildings at {{sample.n_sites}} sites. BDG2 does not state whether timestamps follow the local clock or standard time, so we set each site's mapping from 2016 load ramps at the spring and autumn changes. Four sites are mapped to standard time (Eagle, Moose, Rat, Wolf) and five to clock time with a clear signal (Bear, Bull, Fox, Hog, Peacock). Five ambiguous sites (Bobcat, Cockatoo, Gator, Lamb, Robin) are mapped to clock time, and the opposite mapping is reported as a sensitivity analysis.

Two selection effects need disclosure. First, E1 and E3 read 2017 data: E1 counts 2017 valid days, and E3 uses the 2017 warm-season flat share, including event-day hours. They judge completeness and flatness only; no 2017 load level, error or parameter enters before the freeze. Second, E3 removed most meters at Lamb: of {{sample.excl.Lamb.e1|int}} eligible meters, {{sample.excl.Lamb.e3_flat_meter|int}} were flat, leaving {{sample.excl.Lamb.n_after|int}}, so Lamb's values rest on few buildings. A rerun applying E3 to non-event days only (excluding {{nonevent.sample.n_e3_excluded|int}} meters) covers H1 to H3; matching reruns for H4 and H5 were not run.

The unit of estimation is the building-event and the unit of inference is the site. Every test and interval runs over the {{sample.n_sites}} site-level values and does not extend to a wider population of buildings.

\subsection{Proxy events and controls}
\label{sec:data:events}

BDG2 records no demand-response dispatch, so event days are proxies built from weather and the calendar, not records of dispatch. The primary rule (R1) takes, for each site and year, the 10 hottest non-holiday weekdays from 15 May to 30 September, ranked by daily maximum air temperature; we store ranks up to 40 for the dispatch analysis. Holidays follow U.S. federal holidays \citep{us5usc6103}, bank holidays for England and Wales (Lamb and Robin) \citep{uk_bfda1971} and Irish public holidays (Wolf) \citep{ie_owta1997}. The control rule (R3) draws 10 random non-holiday weekdays from the remaining candidates, 5 times, and averages each building's statistic over the draws. The scored window is 14:00 to 18:00, with 16:00 to 20:00 as a secondary window. R1 yields {{sample.n_events.R1|int}} building-events and R3 yields {{sample.n_events.R3|int}}. A secondary rule (R2-ISO) ranks days by the day-ahead forecast peak of the local balancing authority in EIA-930 \citep{eia930} for seven U.S. sites; it lies outside the confirmatory family.

Hot-weekday load may already contain peak avoidance or utility demand response, so the metered series is the no-curtailment counterfactual by assumption only (Section~\ref{sec:discussion}).

\subsection{Information sets and folds}
\label{sec:data:folds}

Baselines are grouped by the hours of the target day they read (Figure~\ref{fig:design}). I-hist uses the building's own history before the event day. I-morn adds its own event-day hours that end at least one hour before the window. I-donor adds the same-day load of other buildings at the site, plus the building's own event-day hours that end at least four hours before the window.

Buildings are split into two disjoint folds per site, A and B, by a seeded permutation of meter identifiers. The split was frozen before the flat-meter filter and is shared by both years. Targets in one fold use donors from the other. Donors and targets share weather and operating schedules, so site shocks are not independent across them and donor-based methods inherit this dependence.

\subsection{Analysis plan, deviations and audit trail}
\label{sec:data:plan}

We wrote the plan before scoring any 2017 load value and froze it with an MD5 manifest (version 2) covering the plan file, the deviations log, the frozen event dates, holidays and folds, and every pipeline source file \citep{nosek2018}. Five tests form one Holm family \citep{holm1979}. The status line inside the plan file predates the freeze and still reads as a draft; we left the file unedited so that its hash holds, and the dated manifest records the freeze. Two choices were made with development data in view. The absolute-error-growth definition and two-sided test for H5 replaced a signed version already seen on 2016, so that definition is not independent of 2016, and the signed contrast is secondary. The synthetic-control settings (seven pre-period days, ridge penalty of three) \citep{sievers2026} were retuned on 2016 when the folds were frozen.

One deviation from the frozen plan (DEV-1) occurred, and later choices are dated in a post-freeze register (Supplement Section~\ref{sec:S1:dev}). The first 2017 run stopped with a missing-key error for the secondary classic synthetic control in the final-mode code path, before any statistic, ledger row or figure existed; we had logged only meter counts and timings and had opened no event frame. We changed the code only: the final mode now computes both synthetic controls, and a development flag runs the final path on 2016, where it reproduced the reference values exactly. We reissued the manifest as version 2 and kept version 1. A coding error in an early 2016 pilot of the repeated-dispatch analysis, which curtailed donors as well as participants, was corrected before the freeze. The pre-registered LightGBM and multilayer perceptron (MLP) searches ({{hpo.lgbM.trials|int}} trials per variant, tuned on 2016 only) ran after the confirmatory evaluation. Their spaces and budget are pre-registered, and the tuning score was set after the freeze and before any 2017 evaluation of the tuned variants. LightGBM results are reported beside the fixed-parameter values; the MLP is reported only in its tuned form. We cut the augmented synthetic control. The audit trail is in the archive and Supplement Section~\ref{sec:S1}.
